\documentclass{article}
\usepackage[utf8]{inputenc}
\usepackage{geometry}
\geometry{a4paper, margin=1in}
\usepackage{booktabs}
\usepackage{graphicx}
\usepackage{amsmath}
\usepackage{hyperref}

\title{ESC204 Comprehensive Verification Dossier: Systems Integration \& Fluid Dynamics}
\author{ESC204 Engineering Team}
\date{\today}

\begin{document}
\maketitle

\section{System Integration \& Mechanical Verification}
\subsection{The Automated Ridge Track \& Drive Architecture}
The core mechanical verification involved testing the complete integration of the motor-driven ridge track, the custom pulley system, and the dual-nozzle carriage. The primary objective was to empirically verify that the Stepper/Raspberry Pi Pico architecture could autonomously and consistently drive the clearing mechanism across a 60\,cm physical track mock-up without stalling, slipping the drive belt, or overheating under continuous load.

The mechanical drivetrain consists of a Nema 17 Stepper motor connected to a continuous loop pulley system. The carriage is directly affixed to this belt. This system replaced a previously considered direct-drive DC motor approach, which was abandoned due to an inability to perform accurate time-based dead reckoning. By utilizing a stepper architecture, the microcontroller (Raspberry Pi Pico) can command an exact number of steps, translating to millimeter-perfect linear actuation regardless of the variable mass or friction of the carriage.

\subsection{Thermal Management \& Power Segregation}
A critical integration hurdle was managing the power distribution and thermal output of the A4988 stepper driver. Initial integration attempts revealed that pulling high continuous current for the Nema 17 motor would cause logic brownouts if the motor and microcontroller shared a power rail. Consequently, the system was physically segregated: the Raspberry Pi Pico logic operates on an isolated 5V USB source (providing clean 3.3V logic to the A4988 $V_{DD}$ pin), while the motor operates on a dedicated 12V DC wall adapter routed strictly to the $V_{MOT}$ pin.

Furthermore, to ensure the system could run continuously without entering thermal shutdown (a common failure mode for A4988 chips lacking active cooling), the driver's current-limiting potentiometer ($V_{REF}$) was manually tuned. The potentiometer was dialed back to approximately 50\% capacity. This tuning establishes the optimal engineering balance: it provides sufficient amperage to maintain the holding and driving torque required to pull the carriage and attached air hose, while strictly capping the current to prevent the IC from overheating and halting the clearing cycle.

\subsection{Endurance Testing \& Automated Reversal}
The physical track incorporates mechanical limit switches (GP15 and GP16) wired with internal pull-up resistors to the Pico. These switches act as absolute zeroing mechanisms. 

The integrated system was subjected to a rigorous endurance test. The carriage successfully and smoothly traversed the full 60\,cm track. Upon physically striking the extremity limit switches, the CircuitPython logic instantaneously captured the state change (triggering a LOW signal) and executed an immediate reversal of the DIR (direction) pin. 

While the verification protocol only mandated a minimum of 2 complete passes to prove functionality, the system successfully completed a continuous, uninterrupted 2-minute cycle. During this 120-second stress test, the track exhibited zero stalling, the belt maintained tension during sudden direction reversals, and the A4988 driver remained within safe thermal operating limits.

\vspace{0.5cm}
\noindent
\textbf{Video Artifacts (Track Automation):} \href{https://drive.google.com/drive/folders/1wNK_jYUvHXoMne_cIdYbAIURdvYpuG-J?usp=sharing}{View Google Drive Folder}
\begin{itemize}
    \item \textit{File 1:} 8-second sub-clip demonstrating limit switch actuation.
    \item \textit{File 2:} Full 2-minute continuous endurance cycle.
\end{itemize}

\section{Fluid Dynamics \& Clearing Efficiency Verification}
\subsection{Aerodynamic Vector Justification ($30^\circ$ Angle of Attack)}
A fixed angle of attack of $30^\circ$ was selected for the nozzle carriage mount. This decision was justified through a vector decomposition of the airflow. Research into industrial air knives indicates that an angle between $25^\circ$ and $35^\circ$ is mathematically optimal for surface clearing. 

At exactly $30^\circ$, the high-velocity air stream is decomposed into two critical force vectors:
\begin{enumerate}
    \item \textbf{Normal Force ($F_y$):} Sufficient downward pressure to maintain boundary layer attachment (the Coanda effect) and force the air underneath the leading edge of the debris, but not so much pressure that it pins the debris to the roof via static friction.
    \item \textbf{Tangential Shear Stress ($F_x$):} Maximized horizontal force to physically propel the dislodged mass down the slope.
\end{enumerate}
Steeper angles ($>45^\circ$) cause downward pinning (amplifying stiction), while shallower angles ($<15^\circ$) cause the air to deflect harmlessly over the aerodynamic profile of the debris.

\subsection{Experimental Setup: High-Friction Modeling}
To accurately replicate the fluid dynamics of a real-world scenario, a scaled roof mock-up (35\,cm x 15\,cm) was constructed. Crucially, the substrate was covered in 40-grit aluminum oxide sandpaper to accurately model the high coefficient of static friction inherent to asphalt shingles. The mock-up was locked at a $30^\circ$ incline. A custom alignment frame and protractor were utilized to ensure the nozzle outputs were perfectly parallel to the $30^\circ$ pitch. Axial alignment was strictly verified by a secondary observer to eliminate angular bias during data collection. Disaggregated tissue paper was utilized to simulate both dry and heavily saturated autumnal debris.

\subsection{Distance \& Diffusion Trials}
Measurements were taken at three intervals ($d = 10\,$cm, $25\,$cm, and $50\,$cm). Three distinct trials were conducted per configuration to ensure statistical reliability. 
\begin{itemize}
    \item \textbf{Close Range ($d = 10\,$cm):} Both the Custom Air Blade and the standard Circular Nozzle produced sufficient instantaneous impact force to clear the dry debris entirely. Geometrical advantages are marginalized at extreme close range.
    \item \textbf{Medium Range ($d = 25\,$cm):} Fluid dynamic divergence became measurable. The Custom Air Blade ejected the tissue paper significantly higher into the air and further off the mock-up compared to the Circular Nozzle. This empirically proves that the planar shear layer penetrates \textit{underneath} the debris, whereas the circular column merely pushes against its face.
    \item \textbf{Far Range ($d = 50\,$cm):} The Circular Nozzle suffered from severe turbulent diffusion, clearing only a narrow, localized central path and leaving substantial debris along the lateral edges. The Custom Air Blade maintained a coherent, wide-area planar sweep across the entire 15\,cm width. This proves the custom nozzle's superior lateral efficiency, which is critical for minimizing the number of passes the moving carriage must make.
\end{itemize}

\subsection{The "Wet Debris" Maximum Stiction Stress Test}
The most critical verification was the Wet Debris test at $d = 50\,$cm. Saturated tissue paper maximizes static friction and minimizes the aerodynamic profile, removing edges for the air to "grab."
\begin{itemize}
    \item \textbf{Circular Nozzle ($30^\circ$):} Failed. The localized, high-pressure column pushed directly \textit{down} on the low-profile wet tissues, physically pinning them against the 40-grit surface. Only a few central pieces were dislodged.
    \item \textbf{Custom Air Blade ($30^\circ$):} Succeeded. The narrow slit generated a high-velocity shear layer that successfully slipped beneath the water/sandpaper barrier, peeling the saturated debris off the high-friction surface.
\end{itemize}

\vspace{0.5cm}
\noindent
\textbf{Video Artifacts (Air Blade Trials):} \href{https://drive.google.com/drive/folders/1RAJA95EF-XjMrkhjWkuoIBbONH_DbRVu?usp=drive_link}{View Google Drive Folder}
\begin{itemize}
    \item \textit{Contents:} 8 fully documented and labeled final trials (Dry Circular, Dry Air Blade at varying distances, and the definitive Wet Debris stress tests).
\end{itemize}

\end{document}
